\chapter{Amputation Rehabilitation}

\section*{Definition and Overview}
Amputation rehabilitation is the coordinated process of helping a person recover function, independence, and quality of life after the removal of a limb, or part of a limb. The limb may have been lost through injury, surgical treatment of disease, or congenital absence. Rehabilitation ideally begins before surgery with assessment and education, and continues for months to years afterwards. It is delivered by a multidisciplinary team that typically includes doctors, nurses, physiotherapists, occupational therapists, prosthetists, and psychologists. The overall aims are to fit the person for a prosthesis when appropriate, to retrain daily living skills, to manage phantom and stump pain, to prevent secondary complications, and to support emotional adjustment to life with a changed body.

\section*{Epidemiology}
Amputation is relatively uncommon in the general population but is more frequent in older adults and in people with chronic disease. In developed countries, most limb amputations result from complications of peripheral artery disease and diabetes, and the majority involve the lower limb. Amputations from trauma are more common in younger adults and in men. In recent decades, rates of major amputation have fallen in some groups because of better vascular surgery, revascularisation, and diabetic foot care, while advances in prosthetics and rehabilitation services have enabled more people to achieve independent mobility after limb loss. Rates of successful prosthetic use and return to work have consequently improved.

\section*{Aetiology and Pathophysiology}
The underlying causes of amputation shape the rehabilitation pathway:

\begin{itemize}
    \item \textbf{Peripheral vascular disease:} poor blood supply, often from atherosclerosis and diabetes, may make a limb unsalvageable.
    \item \textbf{Diabetes:} neuropathy, infection, and ischaemia combine to cause non-healing foot ulcers that may progress to gangrene.
    \item \textbf{Trauma:} severe crush, blast, or degloving injuries, and burns, may require amputation.
    \item \textbf{Malignancy:} bone and soft tissue tumours sometimes require limb removal to achieve cure.
    \item \textbf{Infection:} severe sepsis or necrotising infection that cannot otherwise be controlled.
    \item \textbf{Congenital limb deficiency:} a limb is absent or underdeveloped from birth.
\end{itemize}

After amputation, the remaining limb undergoes changes in blood flow, muscle bulk, bone, and sensation. The stump heals, shrinks, and matures over weeks to months, and rehabilitation aims to make the most of this altered anatomy while preventing contractures, muscle wasting, and skin problems.

\section*{Clinical Features}
The features of recovery vary with the level of amputation and the person's health:

\begin{itemize}
    \item \textbf{Stump changes:} swelling (oedema), tenderness, and bruising in the early weeks after surgery.
    \item \textbf{Phantom limb sensations:} the feeling that the amputated limb is still present.
    \item \textbf{Phantom limb pain:} painful sensations felt in the missing limb, often described as burning, stabbing, or cramping.
    \item \textbf{Stump pain:} pain localised to the remaining limb, sometimes arising from cut nerve endings, scar tissue, or bone.
    \item \textbf{Reduced mobility and balance:} difficulty walking, transferring, and climbing stairs.
    \item \textbf{Loss of independence:} difficulty with dressing, bathing, and household tasks.
    \item \textbf{Psychological distress:} grief, low mood, body-image concerns, and anxiety about the future.
\end{itemize}

\section*{Differential Diagnosis}
Symptoms during rehabilitation may be confused with other problems:

\begin{itemize}
    \item Stump pain from a neuroma, compared with stump infection, bone spur, or prosthetic problems.
    \item Phantom limb pain, compared with referred pain from the spine or the opposite limb.
    \item Oedema of the stump, compared with deep vein thrombosis or haematoma.
    \item Poor prosthetic fit, compared with poor stump conditioning or weight change.
    \item Low mood as part of normal adjustment, compared with clinical depression requiring treatment.
\end{itemize}

\section*{When to Seek Medical Care}
Seek medical advice promptly during rehabilitation if:

\begin{itemize}
    \item The stump becomes increasingly red, swollen, hot, or discharges pus.
    \item The wound edges separate or the wound breaks open.
    \item There is sudden severe pain, or the stump becomes cool, pale, or numb.
    \item Fever or a general feeling of being unwell develops.
    \item Falls occur, or there is new difficulty with balance or transfers.
    \item Persistent low mood, tearfulness, or thoughts of self-harm arise.
\end{itemize}

\textbf{Call 000 (or local emergency services) for:}
\begin{itemize}
    \item Uncontrolled bleeding from the stump.
    \item Sudden chest pain or breathlessness, which may indicate a blood clot in the lung.
    \item Sudden swelling or pain in the calf or leg, which may indicate a deep vein thrombosis.
    \item Collapse or loss of consciousness.
\end{itemize}

\section*{Diagnosis and Investigations}
Assessment spans the pre-operative period and the months of recovery:

\begin{itemize}
    \item \textbf{Pre-amputation assessment:} physical fitness, vascular status, cognitive function, and functional goals, so that rehabilitation can be planned before surgery.
    \item \textbf{Wound and stump review:} regular inspection for healing, infection, and pressure areas.
    \item \textbf{Imaging:} X-ray, ultrasound, or CT/MRI to assess bone, blood vessels, or suspected complications.
    \item \textbf{Vascular studies:} to confirm blood supply where limb salvage is being considered.
    \item \textbf{Pain assessment:} questionnaires and examination to distinguish phantom from stump pain and to monitor treatment.
    \item \textbf{Functional assessment:} strength, balance, and activities of daily living measured by physiotherapists and occupational therapists.
\end{itemize}

\section*{Management}
Rehabilitation is delivered by a multidisciplinary team in phases:

\begin{itemize}
    \item \textbf{Prehabilitation:} conditioning, education, and psychological preparation before planned surgery.
    \item \textbf{Stump care:} wound dressing, compression or shaping of the stump, and careful skin hygiene.
    \item \textbf{Pain management:} medicines, nerve blocks, mirror therapy, and psychological techniques for phantom and stump pain.
    \item \textbf{Physiotherapy:} range-of-motion exercises, strengthening, balance training, and gait re-education.
    \item \textbf{Occupational therapy:} retraining for dressing, bathing, cooking, and use of assistive equipment, with home modifications as needed.
    \item \textbf{Prosthetics:} socket fitting, alignment, and gait training with a prosthesis where suitable.
    \item \textbf{Psychological support:} counselling, peer support, and treatment of depression or anxiety.
    \item \textbf{Follow-up:} regular prosthetic review, monitoring of the opposite limb, and lifelong skin and vascular checks.
\end{itemize}

\section*{Complications}
Complications can arise during rehabilitation and require prompt recognition:

\begin{itemize}
    \item Wound infection, breakdown, or delayed healing of the stump.
    \item Haematoma or seroma formation.
    \item Phantom limb pain and chronic stump pain.
    \item Neuroma formation at the cut nerve end.
    \item Joint contractures and muscle wasting if the limb is not moved appropriately.
    \item Skin problems from prosthetic wear, including pressure sores, sweating, and irritation.
    \item Deep vein thrombosis and falls, especially in the early period.
    \item Depression, anxiety, and social isolation.
    \item Increased workload on the remaining limb, predisposing to overuse injury.
\end{itemize}

\section*{Prognosis and Outcomes}
Most people with a lower-limb amputation achieve good mobility, particularly if they were active before surgery and receive early rehabilitation. Many upper-limb amputees learn to use a functional or cosmetic prosthesis, although some prefer to manage one-handed. Outcomes are better when rehabilitation begins promptly and the person is supported by family and an experienced team. Older adults and those with vascular disease have higher rates of re-amputation and reduced life expectancy due to their underlying illness, but most regain the ability to live independently. Ongoing adjustment is normal, and many people report a good quality of life after an initial period of loss and adaptation.

\section*{Prevention and Self-Care}
Self-care protects the remaining limbs and supports recovery:

\begin{itemize}
    \item Manage diabetes, blood pressure, and cholesterol carefully to protect remaining limbs.
    \item Check both feet daily and attend regular podiatry if you have diabetes or poor circulation.
    \item Stop smoking, which damages blood vessels and impairs wound healing.
    \item Exercise within rehabilitation limits to maintain strength, balance, and cardiovascular fitness.
    \item Care for the stump daily: keep it clean and dry, and inspect it for pressure marks.
    \item Wear a well-fitted prosthesis and attend regular prosthetic appointments.
    \item Follow a balanced diet with enough protein to support healing.
    \item Seek psychological support early if grief or low mood is affecting recovery.
\end{itemize}

\section*{Sources and Further Reading}
The following reputable sources provide reliable, up-to-date information on amputation rehabilitation:

\begin{itemize}
    \item NHS. \textit{Amputation and rehabilitation information.} https://www.nhs.uk/
    \item Mayo Clinic. \textit{Prosthetics and rehabilitation services information.} https://www.mayoclinic.org/
    \item MSD Manual. \textit{Rehabilitation after amputation overview.} https://www.msdmanuals.com/
    \item MedlinePlus. \textit{Amputation and prosthetic rehabilitation information.} https://medlineplus.gov/
    \item World Health Organization. \textit{Information on rehabilitation and assistive technology.} https://www.who.int/
    \item NICE. \textit{Clinical guidance on rehabilitation pathways.} https://www.nice.org.uk/
\end{itemize}
